\section{Appendix: Scope and Qualifications}
\vspace{-0.6em}
\begingroup
\setlength{\parskip}{0.3em}%

\textbf{Edition and formulary identity.} This guide covers one formulary: the
Mass of the Eighth Sunday after Pentecost, \latin{Dominica Octava post
Pentecosten}, \latin{II classis}, as printed at pages 387--388 of the
\work{Missale Romanum}, editio typica (Vatican City: Typis Polyglottis
Vaticanis, 1962), marginal numbers 1512--1521. The rank and its occurrence
effects come from the same book's \latin{Rubricæ generales}, nn.\ 10--18,
printed page XII. In the repository's stable temporal series this leaf is
catalogue identity \texttt{22}, a catalogue position and not an occurrence
schedule; the Sunday fell on 19 July in 2026, but the guide is dated to no civil
year.

\textbf{Text verification.} Every appointed Latin form, heading, rank, marginal
number, scriptural reference, ligature, accent, and formulary boundary was
collated on 2026-07-25 against 300-dpi page images of the Church Music
Association of America facsimile of the 1962 Vatican typical edition, SHA-256
\texttt{648fdb8f\dots3518a}, read column by column. The facsimile's embedded
machine text located the leaf only and is preserved unedited in
\texttt{propers/retrieved.txt}; no second image witness was needed, because no
reading was unclear. Transcriptions and recorded divergences are in
\texttt{propers/verified.md}.

\textbf{Translation and rights.} The 1962 Latin liturgical text is public domain
and the project asserts no rights in it. Every English rendering of an appointed
text here is quoted from one identified public-domain witness --- \work{The Roman
Missal translated into the English language for the use of the laity}
(Philadelphia: Eugene Cummiskey, 1861), printed pp.\ 411--413, read as page
images in the Internet Archive scan \texttt{romanmissaltran00churgoog}. That is a
nineteenth-century hand-missal translation with pre-1962 headings and old
orthography: not an approved liturgical translation, quoted as a historical
witness, and nothing here may be used for recitation. Scripture quoted outside
the appointed texts is Douay--Rheims (Challoner) and Clementine Vulgate at
\texttt{drbo.org}; patristic quotations in English are NPNF, and Cyril of
Alexandria is R.~Payne Smith's translation from the Syriac. Neither the project
licence nor this guide relicenses Scripture, liturgical text, or any third-party
translation.

\textbf{Source scope, and what it does not cover.} Patristic, medieval, and
Doctoral reception was searched for every directly appointed passage across the
corpora named in \texttt{research/scope.md}: New Advent, CCEL's NPNF volumes,
Documenta Catholica Omnia, the \work{Patrologia Latina} text at Corpus Corporum,
Corpus Thomisticum, and the Internet Archive. Every retained witness was read at
its own work and locus. The search is bounded and correctable: no critical
edition (CSEL, CCSL, SC, Leonine) was used; machine-read Migne is flagged at the
citation; Greek witnesses were reached in translation except where a variant is
discussed; and no claim is made that all witnesses everywhere have been found.

\textbf{Material negative results.} Ambrose does not comment on the narrative of
Luke 16:1--8 at all; his treatment covers verses 13, 9 and 12, in that order.
Gregory the Great has no homily on this pericope --- his Luke 16 homily is on
Dives and Lazarus --- so his contribution comes from the \work{Moralia}. No
connected Chrysostom exposition of Lk 16:1--9 was located, and Jerome records
that he searched for Origen's and Didymus's explanations and found neither.
These absences are evidence, recorded rather than papered over.

\textbf{Disagreement retained.} On the Gospel this guide preserves five
unresolved cruces: who commends, whether the fraud is censured, what makes the
mammon iniquitous, who receives into the eternal dwellings, and whether it is the
hearer or the money that fails. No consensus is asserted, and agreement should
not be inferred from the proximity of names.

\textbf{Unresolved textual question.} Psalm 47:2 is printed twice in this one
formulary in two Latin forms, \latin{laudábilis nimis} at the Introit and
\latin{laudábilis valde} at the Alleluia. The divergence is transmitted, not
peculiar to this edition; assigning the two forms to named Latin psalters would
need collation against psalter witnesses this study did not undertake, and the
question is left open.

\textbf{Claim classes.} Textual observation, documented historical orientation,
documented reception, and source-grounded synthesis appear together in the map,
the thematic section, and the commentary. Exploratory editorial and AI proposal
appears only in \emph{Interpretive Possibilities}, which carries its own notice;
nothing there is attributed to any source.

\textbf{Review state.} This guide is generation-provenanced and source-audited
by its own records. It has received no independent specialist, theological,
canonical, or ecclesiastical review, and none is claimed. It is a hand-missal
companion for study and prayerful understanding: not a liturgical book, not a
critical edition, not a homily, not a substitute for its sources.

\textbf{As of, and audit records.} Checks and collations reported here were made
on 2026-07-25 unless another date is given at the citation; web witnesses are
mutable, and only the controlling facsimile is pinned by hash. Retrieval
mechanics, the reception matrix, the notable-and-quotable and
interpretive-proposal audits, rejected leads, and negative results are in
\texttt{research/scope.md} and \texttt{research/source-audit.md}; verified texts
in \texttt{propers/verified.md}; bound sources in
\texttt{research/source-bindings.toml}.

\endgroup

\clearpage
